\subsection{\texorpdfstring{第~\ref{sec:holonomic-constraints}~节}{完整约束}}\label{subsec:solutions-holonomic-constraints}

\begin{description}[font=\normalfont,leftmargin=2em,labelwidth=1.5em]
\RGsolution{ex:double-pendulum-motion}{double-pendulum-motion}
\begin{enumerate}[label=(\alph*),ref=\alph*,leftmargin=*]
\item 粒子 $m_1$ 的位置是
\[
  (x_1,y_1)=l_1(\sin\theta,-\cos\theta).
\]
所以其速度为
\[
  (\dot{x}_1,\dot y_1)=l_1\dot\theta(\cos\theta,\sin\theta),
\]
从而动能为
\[
  K_1=\frac12m_1l_1^2\dot\theta^2.
\]
类似地, 粒子 $m_2$ 的位置为
\[
  (x_2,y_2)=l_1(\sin\theta,-\cos\theta)+l_2(\sin\varphi,-\cos\varphi),
\]
其速度为
\[
  (\dot{x}_2,\dot y_2)=l_1\dot\theta(\cos\theta,\sin\theta)+l_2\dot\varphi(\cos\varphi,\sin\varphi),
\]
对应的动能为
\[
  K_2=\frac12m_2l_1^2\dot\theta^2+\frac12m_2l_2^2\dot\varphi^2
    +m_2l_1l_2\cos(\theta-\varphi)\dot\theta\dot\varphi.
\]
于是动能映射 $K:TN\to\mathbb{R}$ 在局部坐标中由
\[
  K=\frac12(m_1+m_2)l_1^2(v^\theta)^2+\frac12m_2l_2^2(v^\varphi)^2
    +m_2l_1l_2\cos(\theta-\varphi)v^\theta v^\varphi
\]
给出. 因为势能显然为
\[
  U=m_1gy_1+m_2gy_2=-(m_1+m_2)gl_1\cos\theta-m_2gl_2\cos\varphi,
\]
所以运动方程为
\begin{align*}
  &(m_1+m_2)l_1^2\ddot\theta+m_2l_1l_2\cos(\theta-\varphi)\ddot\varphi\\
  &\qquad+m_2l_1l_2\sin(\theta-\varphi)\dot\varphi^2=-(m_1+m_2)gl_1\sin\theta
\end{align*}
与
\begin{align*}
  &m_2l_2^2\ddot\varphi+m_2l_1l_2\cos(\theta-\varphi)\ddot\theta\\
  &\qquad-m_2l_1l_2\sin(\theta-\varphi)\dot\theta^2=-m_2gl_2\sin\varphi.
\end{align*}
\item 线性化方程为
\[
  (m_1+m_2)l_1^2\ddot\theta+m_2l_1l_2\ddot\varphi=-(m_1+m_2)gl_1\theta
\]
与
\[
  m_2l_2^2\ddot\varphi+m_2l_1l_2\ddot\theta=-m_2gl_2\varphi.
\]
这些方程中满足 $\varphi=k\theta$ 的解必须同时满足
\[
  (m_1l_1^2+m_2l_1^2+km_2l_1l_2)\ddot\theta=-(m_1+m_2)gl_1\theta
\]
和
\[
  (km_2l_2^2+m_2l_1l_2)\ddot\theta=-km_2gl_2\theta.
\]
因此 $k$ 必须满足
\[
  \begin{vmatrix}
    (m_1+m_2)l_1^2+km_2l_1l_2&(m_1+m_2)gl_1\\
    km_2l_2^2+m_2l_1l_2&km_2gl_2
  \end{vmatrix}=0,
\]
即
\[
  k=\frac{l_2-l_1\pm\sqrt{(l_2-l_1)^2+4\mu l_1l_2}}{2\mu l_2},
\]
其中
\[
  \mu=\frac{m_2}{m_1+m_2}\in(0,1).
\]
注意 $k$ 有两个可能值, 一正一负, 分别对应两摆同相或反相振荡.

由第一个线性化运动方程可知, 振荡周期为
\[
  2\pi\sqrt{\frac{l_1+k\mu l_2}{g}}
  =2\pi\sqrt{\frac{l_1+l_2\pm\sqrt{(l_2-l_1)^2+4\mu l_1l_2}}{2g}}.
\]
注意, 两摆同相振荡时周期较长, 反相振荡时周期较短.
\end{enumerate}
\end{description}
